% !TeX root = ../../../金融经济学_王江_讲义及习题.tex
% 示意取 c=d=1、a/c=2，均值偏差分别为 1.25 和 -.8，乘积为 -1。
\begin{tikzpicture}[
 >=Stealth,font=\small,
 frontier/.style={blue!70!black},
 zero/.style={teal!65!black},
 guide/.style={gray,densely dotted}]
 \begin{scope}[x=1.50cm,y=.96cm]
  \coordinate (P) at ({sqrt(1+1.25^2)},3.25);
  \coordinate (MVP) at (1,2);
  \coordinate (ZCP) at ({sqrt(1+.8^2)},1.20);
  \draw[->] (0,.46)--(2.57,.46) node[right] {$\sigma$};
  \draw[->] (0,.46)--(0,3.88) node[above] {$\bar r$};
  \draw[guide] (0,3.25) node[left,black] {$\bar r_p$}--(P);
  \draw[guide] (0,2) node[left,black] {$a/c$}--(MVP);
  \draw[guide] (0,1.20) node[left,black] {$\bar r_{zcp}$}--(ZCP);
  \draw[frontier,very thick,domain=-1.32:1.64,samples=81,variable=\t]
    plot ({sqrt(1+\t*\t)},{2+\t});
  \fill[frontier] (P) circle (2pt) node[right=3pt] {$p$};
  \fill[black] (MVP) circle (1.8pt) node[right=3pt] {MVP};
  \fill[zero] (ZCP) circle (2pt) node[right=3pt] {$zcp$};
  \node at (1.22,-.32) {(a) 前沿上的两个组合};
 \end{scope}
 \begin{scope}[xshift=5.85cm,x=1.45cm,y=1.34cm]
  % 横坐标按 r_q-a/c 加上 1.25 平移，协方差的零水平设为 .75。
  \coordinate (root) at (.45,.75);
  \coordinate (mvpCov) at (1.25,1.75);
  \draw[->] (0,.75)--(3.25,.75) node[right] {$\bar r_q$};
  \draw[->] (0,.28)--(0,2.76) node[above right=-1pt] {$\operatorname{Cov}(\widetilde r_p,\widetilde r_q)$};
  \node[left] at (0,.75) {$0$};
  \draw[guide] (0,1.75) node[left,black] {$1/c$}--(mvpCov)--(1.25,.75)
    node[below,black] {$a/c$};
  \draw[frontier,very thick] (.10,.3125)--(2.02,2.7125);
  \node[frontier,anchor=west,inner sep=0pt] at (2.10,2.48)
    {$\begin{aligned}
      \operatorname{Cov}(\widetilde r_p,\widetilde r_q)
        &=\frac{b-a\bar r_p}{d}+\frac{c\bar r_p-a}{d}\,\bar r_q\\
        &=\frac{c\bar r_p-a}{d}\,(\bar r_q-\bar r_{zcp}).
    \end{aligned}$};
  \fill[black] (mvpCov) circle (1.8pt);
  \fill[zero] (root) circle (2pt) node[below right=2pt] {$\bar r_{zcp}$};
  \node at (1.55,{-.32*.96/1.34}) {(b) 协方差随 $q$ 的均值变化};
 \end{scope}
\end{tikzpicture}
